\documentclass{amsart}
\usepackage{amsmath,amssymb,amsthm,hyperref}

\newtheorem{theorem}{Theorem}
\newtheorem{lemma}{Lemma}
\newtheorem{proposition}{Proposition}

\begin{document}

\title{No nontrivial irreducible finite $k$-scheme is rigid}
\maketitle

\begin{theorem}\label{thm:main}
Let $k$ be an algebraically closed field. There is \emph{no} irreducible finite
$k$-scheme $\Gamma$ other than $\operatorname{Spec}(k)$ with $T^1_\Gamma=0$.
Equivalently, every finite-dimensional local commutative $k$-algebra $A\neq k$
admits a nontrivial first-order infinitesimal deformation; that is, $T^1_A\neq 0$.
\end{theorem}

Throughout, $A$ denotes a finite-dimensional commutative local $k$-algebra with
maximal ideal $\mathfrak m$ and residue field $A/\mathfrak m=k$ (the residue field
is $k$ because $A$ is finite over the algebraically closed field $k$). All
algebras are commutative, associative, unital $k$-algebras. We assume $A\neq k$,
equivalently $\mathfrak m\neq 0$.

\section{The Lichtenbaum--Schlessinger description of $T^1$}

We recall the cotangent-complex (Lichtenbaum--Schlessinger) computation of the
first-order deformation module of a $k$-algebra, in the form we need. This is the
standard, characteristic-free model for $T^1$.

\subsection{A presentation of $A$}
Set $n=\dim_k \mathfrak m/\mathfrak m^2$, the embedding dimension. Choose
$y_1,\dots,y_n\in\mathfrak m$ whose images form a $k$-basis of
$\mathfrak m/\mathfrak m^2$. Let
\[
R=k[x_1,\dots,x_n]_{(x_1,\dots,x_n)}
\]
be the localization of the polynomial ring at its irrelevant maximal ideal; $R$
is a Noetherian \emph{regular local} ring with maximal ideal
$\mathfrak m_R=(x_1,\dots,x_n)$ and residue field $k$. Define a local
$k$-algebra homomorphism
\[
\pi\colon R\longrightarrow A,\qquad x_i\longmapsto y_i .
\]
Since the $y_i$ generate $\mathfrak m/\mathfrak m^2$ and $\mathfrak m$ is
nilpotent (as $A$ is Artinian local), Nakayama's lemma shows the $y_i$ generate
$\mathfrak m$, hence $\pi$ is surjective. Put $I=\ker\pi$, so $A=R/I$, and $I$ is
a finitely generated ideal of $R$ because $R$ is Noetherian.

\begin{lemma}\label{lem:I-in-msq}
$I\subseteq \mathfrak m_R^{\,2}$ and $I\neq 0$.
\end{lemma}

\begin{proof}
The induced map $\mathfrak m_R/\mathfrak m_R^2\to \mathfrak m/\mathfrak m^2$ sends
$x_i\mapsto y_i$; by the choice of the $y_i$ it is an isomorphism of
$k$-vector spaces. If $f\in I$ then $f\in\mathfrak m_R$ (as $\pi(f)=0\in\mathfrak m$
and $\pi^{-1}(\mathfrak m)=\mathfrak m_R$), and the image of $f$ in
$\mathfrak m_R/\mathfrak m_R^2$ maps to $\pi(f)\bmod\mathfrak m^2=0$; by injectivity
of the displayed map, $f\in\mathfrak m_R^2$. Thus $I\subseteq\mathfrak m_R^2$.

If $I=0$ then $A\cong R$, but $R$ is an integral domain of Krull dimension $n\ge 1$
(note $n\ge 1$ since $\mathfrak m\neq 0$), while $A$ is finite-dimensional over $k$,
a contradiction. Hence $I\neq 0$.
\end{proof}

\subsection{The conormal module and the formula for $T^1$}
Let $I/I^2$ be the conormal $A$-module and let $\Omega_{R/k}$ be the module of
K\"ahler differentials of $R$ over $k$; since $R$ is a localization of a polynomial
ring, $\Omega_{R/k}$ is free with basis $dx_1,\dots,dx_n$, so
$\Omega_{R/k}\otimes_R A\cong A^n$. The truncated cotangent complex of $A/k$ is
\begin{equation}\label{eq:cotangent}
I/I^2 \xrightarrow{\;\;\delta\;\;} \Omega_{R/k}\otimes_R A\cong A^n,
\qquad \delta([f])=\sum_{i=1}^n \pi\!\Big(\tfrac{\partial f}{\partial x_i}\Big)\,
e_i ,
\end{equation}
where $e_i$ is the image of $dx_i$ and $[f]$ denotes the class of $f\in I$ in
$I/I^2$. (The map $\delta$ is well defined: for $f,g\in I$ the Leibniz rule gives
$d(fg)=f\,dg+g\,df$, whose image in $\Omega_{R/k}\otimes_R A$ vanishes because
$\pi(f)=\pi(g)=0$; hence $\delta(I^2)=0$.)

Applying $\operatorname{Hom}_A(-,A)$ to \eqref{eq:cotangent} yields the
cochain complex
\begin{equation}\label{eq:dual-complex}
\operatorname{Hom}_A(A^n,A)\;=\;A^n
\xrightarrow{\;\;J\;\;}
\operatorname{Hom}_A(I/I^2,A),
\qquad
J(\theta)([f])=\sum_{i=1}^n \pi\!\Big(\tfrac{\partial f}{\partial x_i}\Big)\,\theta_i ,
\end{equation}
for $\theta=(\theta_1,\dots,\theta_n)\in A^n$ (we identify
$\theta\in\operatorname{Hom}_A(A^n,A)$ with the tuple $\theta_i=\theta(e_i)$).
Here $J=\operatorname{Hom}_A(\delta,A)$ is the dual of the Jacobian map. The
Lichtenbaum--Schlessinger theory gives a canonical isomorphism
\begin{equation}\label{eq:LS}
T^1_A\;=\;T^1(A/k,A)\;\cong\;\operatorname{coker}\bigl(J\bigr)
=\operatorname{Hom}_A(I/I^2,A)\big/\operatorname{im}(J).
\end{equation}
Indeed, this is the standard model: a first-order deformation of $A$ over
$k[\varepsilon]/(\varepsilon^2)$ amounts to lifting the relations, i.e.\ to an
$A$-linear map $I/I^2\to A$ recording the first-order change of the relations,
and two such lifts give isomorphic deformations precisely when they differ by a
lift coming from an infinitesimal change of the generators $x_i\mapsto x_i+
\varepsilon\theta_i$, which is exactly an element of $\operatorname{im}(J)$. The
identification of this cokernel with the tangent space to the deformation functor
is the content of the Lichtenbaum--Schlessinger construction of the cotangent
complex (and agrees with $T^1(A/k,A)$ for any presentation, in every
characteristic).

\section{Nonvanishing of $T^1$}

We now show $\operatorname{coker}(J)\neq 0$, which by \eqref{eq:LS} proves
Theorem~\ref{thm:main}. The idea is that, because the relations lie in
$\mathfrak m_R^2$ (Lemma~\ref{lem:I-in-msq}), the dual Jacobian $J$ has image
inside maps with values in $\mathfrak m$; reducing the \emph{target} of a homomorphism
modulo $\mathfrak m$ therefore kills $\operatorname{im}(J)$ and detects a nonzero
quotient.

\subsection{A reduction-mod-$\mathfrak m$ map}
Let $q\colon A\twoheadrightarrow A/\mathfrak m=k$ be the quotient map. Define
\begin{equation}\label{eq:rho}
\rho\colon \operatorname{Hom}_A(I/I^2,A)\longrightarrow
\operatorname{Hom}_A(I/I^2,k),\qquad \rho(\varphi)=q\circ\varphi .
\end{equation}
Because $k=A/\mathfrak m$ is annihilated by $\mathfrak m$, every $A$-linear map
$I/I^2\to k$ factors through
$(I/I^2)\big/\mathfrak m\,(I/I^2)$. Since $I^2\subseteq \mathfrak m_R I$ (as
$I\subseteq\mathfrak m_R$) and $\mathfrak m\,(I/I^2)=(\mathfrak m_R I+I^2)/I^2
=\mathfrak m_R I/I^2$, we get
\[
(I/I^2)\big/\mathfrak m\,(I/I^2)\;=\;I/\mathfrak m_R I ,
\]
a finite-dimensional $k$-vector space. Hence $\operatorname{Hom}_A(I/I^2,k)$ is
canonically the $k$-linear dual:
\begin{equation}\label{eq:dual-iso}
\operatorname{Hom}_A(I/I^2,k)\;\cong\;\operatorname{Hom}_k\!\bigl(I/\mathfrak m_R I,\;k\bigr)
=\bigl(I/\mathfrak m_R I\bigr)^{*}.
\end{equation}

\begin{lemma}\label{lem:rho-surj}
The map $\rho$ of \eqref{eq:rho} is surjective.
\end{lemma}

\begin{proof}
Let $\psi\in\operatorname{Hom}_A(I/I^2,k)$. By \eqref{eq:dual-iso}, $\psi$ is a
$k$-linear functional on the $k$-vector space $V:=I/\mathfrak m_R I$. Choose
$f_1,\dots,f_r\in I$ whose classes form a $k$-basis of $V$; by Nakayama's lemma
(over the local ring $R$, applied to the finitely generated $R$-module $I$ with
$I\subseteq\mathfrak m_R$), these $f_j$ generate $I$ as an ideal, hence their
classes $[f_1],\dots,[f_r]$ generate $I/I^2$ as an $A$-module. Define
$\varphi\colon I/I^2\to A$ on these generators by
\[
\varphi([f_j])=\psi(f_j\bmod\mathfrak m_R I)\cdot 1_A\in k\cdot 1_A\subseteq A,
\qquad j=1,\dots,r,
\]
and extend $A$-linearly. This is well defined: any $A$-linear relation
$\sum_j a_j[f_j]=0$ in $I/I^2$ means $\sum_j \tilde a_j f_j\in I^2\subseteq
\mathfrak m_R I$ for lifts $\tilde a_j\in R$ of $a_j\in A$; reducing modulo
$\mathfrak m_R$ gives $\sum_j \bar a_j\,(f_j\bmod\mathfrak m_R I)=0$ in $V$, where
$\bar a_j\in k$ is the residue of $a_j$. Applying the $k$-linear $\psi$ yields
$\sum_j \bar a_j\,\psi(f_j\bmod\mathfrak m_R I)=0$; and since the
$\varphi([f_j])\in k\cdot1_A$ are killed by $\mathfrak m$, we have
$a_j\varphi([f_j])=\bar a_j\,\varphi([f_j])$, so
$\sum_j a_j\varphi([f_j])=\bigl(\sum_j \bar a_j\,\psi(f_j\bmod\mathfrak m_R I)\bigr)1_A=0$.
Thus $\varphi$ respects all relations and is a well-defined element of
$\operatorname{Hom}_A(I/I^2,A)$. By construction $\rho(\varphi)=q\circ\varphi=\psi$,
since $q$ is the identity on $k\cdot1_A$. Hence $\rho$ is surjective.
\end{proof}

\subsection{$\rho$ annihilates $\operatorname{im}(J)$}
\begin{lemma}\label{lem:J-killed}
$\rho\circ J=0$, i.e.\ $\operatorname{im}(J)\subseteq\ker\rho$.
\end{lemma}

\begin{proof}
Let $\theta\in A^n$ and $f\in I$. By Lemma~\ref{lem:I-in-msq},
$f\in\mathfrak m_R^2$, so each partial derivative
$\partial f/\partial x_i\in\mathfrak m_R$; therefore its image
$\pi(\partial f/\partial x_i)\in\mathfrak m$. From \eqref{eq:dual-complex},
\[
J(\theta)([f])=\sum_{i=1}^n
\pi\!\Big(\tfrac{\partial f}{\partial x_i}\Big)\,\theta_i\in\mathfrak m\cdot A=\mathfrak m .
\]
Hence $q\bigl(J(\theta)([f])\bigr)=0$ for all $f$, i.e.\ $\rho(J(\theta))=0$.
\end{proof}

\subsection{Conclusion}
By Lemma~\ref{lem:J-killed}, $\rho$ factors through the cokernel of $J$:
there is a $k$-linear surjection
\[
\bar\rho\colon \operatorname{coker}(J)\;\xrightarrow{\;\;}\;
\operatorname{Hom}_A(I/I^2,k)\cong\bigl(I/\mathfrak m_R I\bigr)^{*},
\]
which is surjective because $\rho$ is (Lemma~\ref{lem:rho-surj}) and
$\operatorname{coker}(J)\to\operatorname{Hom}_A(I/I^2,k)$ is induced by the surjective
$\rho$. Now $I/\mathfrak m_R I\neq 0$: indeed $I\neq 0$ by Lemma~\ref{lem:I-in-msq},
and Nakayama's lemma (over the local ring $R$, $I$ finitely generated,
$I\subseteq\mathfrak m_R$) gives $I=\mathfrak m_R I\Rightarrow I=0$, so $I\neq 0$
forces $I\neq\mathfrak m_R I$, i.e.\ $I/\mathfrak m_R I\neq 0$. Therefore its dual
$(I/\mathfrak m_R I)^{*}\neq 0$, and by the surjection $\bar\rho$,
\[
T^1_A\cong\operatorname{coker}(J)\neq 0 .
\]
This proves Theorem~\ref{thm:main}: every finite-dimensional local commutative
$k$-algebra $A\neq k$ has $T^1_A\neq 0$, so the only irreducible finite
$k$-scheme that is rigid is $\operatorname{Spec}(k)$.
\hfill$\qed$

\section{Remark: an explicit nonzero deformation}

The proof produces an explicit first-order deformation. Choose
$f_1,\dots,f_r\in I$ as in Lemma~\ref{lem:rho-surj} (a minimal system of
generators of $I$, i.e.\ a $k$-basis of $I/\mathfrak m_R I$). Pick any nonzero
functional $\psi\in(I/\mathfrak m_R I)^{*}$, e.g.\ $\psi(f_1\bmod\mathfrak m_R I)=1$
and $\psi(f_j\bmod\mathfrak m_R I)=0$ for $j\ge 2$. The associated deformation of
$A=R/I$ over $k[\varepsilon]/(\varepsilon^2)$ is
\[
A_\varepsilon \;=\; R[\varepsilon]\big/\bigl(\varepsilon^2,\;
f_1-\varepsilon,\;f_2,\dots,f_r\bigr),
\]
i.e.\ the relation $f_1=0$ is deformed to $f_1=\varepsilon$. The class of
$A_\varepsilon$ in $T^1_A$ is $\bar\rho^{-1}$-related to $\psi\neq 0$, hence
nontrivial: $A_\varepsilon$ is a flat $k[\varepsilon]/(\varepsilon^2)$-algebra
that is not isomorphic, as a deformation, to the trivial one
$A\otimes_k k[\varepsilon]/(\varepsilon^2)$. For instance, for
$A=k[x]/(x^2)$ one has $R=k[x]_{(x)}$, $I=(x^2)$, and the deformation is
$k[x,\varepsilon]/(\varepsilon^2,\,x^2-\varepsilon)$, which is the (flat,
nontrivial) first-order smoothing of the double point.

\end{document}
